%
%  chapter1.tex
%
\chapter{Introduction}
\label{chap:introduction}

\section{Context and Motivation}

Education systems worldwide face a persistent challenge: sustaining learner engagement over time, particularly for subjects such as sustainability and digital literacy that require continuous effort rather than a single study session. Unlike topics learned in a short burst, these disciplines ask learners to build habits and update knowledge regularly. Traditional instructional platforms respond to this challenge with static question banks authored by human experts. The problem is that manual authoring is expensive: producing one high-quality educational task may take an experienced author between thirty minutes and several hours. As a platform grows to cover more subjects, more difficulty levels, or more languages, the authoring workforce must grow proportionally. The economics do not improve with scale.

Gamified platforms have emerged as a structural response to the engagement problem. Research consistently shows that game mechanics --- points, streaks, leaderboards, and immediate feedback --- produce measurable improvements in learner persistence and knowledge retention \cite{hamari2014does}. A player who earns points for completing today's task has a reason to return tomorrow. A player whose streak is at risk of a penalty has a reason not to skip a day. These mechanisms convert educational participation from an act of willpower into a structured behaviour reinforced by the platform itself.

However, two structural problems remain unsolved in combination by existing platforms. First, the content authoring bottleneck: even the most successful gamified platforms, such as Duolingo, depend on large teams of human authors to produce and maintain content. Automated content generation using AI models offers a path to removing this bottleneck, but doing so requires careful engineering --- AI inference is slow, expensive, and non-deterministic, and cannot simply be placed on the path of a user's request. Second, the mobile platform fragmentation problem: the two dominant mobile ecosystems, Android and iOS, use different programming languages and toolchains. A team that builds and maintains two separate codebases doubles engineering cost, doubles the time required to ship any feature, and doubles the surface area for bugs. Cross-platform development frameworks have matured considerably in recent years, but their systematic adoption in gamified educational platforms remains limited.

U!REKA Play4Change was designed to address both problems simultaneously. The platform produces educational tasks automatically, using a scheduled AI content generation pipeline that runs nightly and stores results before users ever request them. It delivers an Android application and an iOS application from a single shared codebase, using KMP (Kotlin Multiplatform) --- a framework that allows the same business logic and user interface code to run on both platforms. The result is a platform that scales content without scaling authoring effort, and scales across mobile platforms without scaling the engineering team.

\section{Problem Statement}

Existing gamified educational platforms cannot scale content production without proportionally scaling human authoring effort. The platforms that have succeeded in sustaining engagement over time --- Duolingo being the clearest example --- do so through the work of large, specialised content teams. This model is not accessible to small educational institutions, individual researchers, or non-profit organisations with limited resources. Automated content generation using large language models (LLMs --- AI systems trained to produce coherent text from instructions) offers a path to removing this constraint, but introduces its own engineering challenges: inference latency, availability dependency on third-party AI services, cost that scales with users, and the risk of generating repetitive content over time.

Simultaneously, building native applications for both Android and iOS requires separate development teams, separate codebases, and separate testing and maintenance cycles. For organisations without the resources of large technology companies, this fragmentation is a practical barrier to multiplatform deployment. Current cross-platform solutions each make different trade-offs between code sharing, performance, and platform fidelity. None has been adopted by a production gamified educational platform that also solves the AI content generation problem.

No existing platform solves both problems in combination. U!REKA Play4Change addresses this gap.

\section{Proposed Solution}

U!REKA Play4Change is a gamified educational platform built on four interconnected engineering components.

The first is a single Kotlin Multiplatform codebase delivering both Android and iOS applications. Shared business logic --- the scoring engine, data models, network clients, and local persistence --- is written once and compiled for both platforms. The shared declarative user interface, built with Compose Multiplatform, renders natively on each platform without visual compromise.

The second is a nightly AI content generation pipeline. A scheduled batch job runs each night, invoking the Mistral language model API through the LangChain4j library to generate the educational tasks for the following day. By separating the moment of generation from the moment of delivery, the platform avoids placing AI inference latency on the user's request path. Content is ready in storage when users arrive.

The third is semantic deduplication using pgvector (a vector similarity search extension for the PostgreSQL database). Because language models produce different phrasings of the same concept across separate invocations, exact text comparison cannot detect repetition. Each generated task is converted to a numerical vector --- a representation of its meaning rather than its words --- and compared against all previously stored vectors. Tasks that are conceptually too similar to existing content are discarded before storage.

The fourth is a server-side scoring engine that awards points for correct and timely task completions, applies penalties for inactivity, and manages badge awards. The scoring logic lives on the server, not the client, to prevent manipulation and to ensure correctness under concurrent submissions.

The backend runs on Spring Boot (a Java framework for building server applications) and is deployed as a set of coordinated containers using Docker Compose (a tool that manages multiple application services as a single unit).

\section{Objectives}

The project pursued four core engineering objectives.

\textbf{Objective 1: Multiplatform delivery from a single codebase.} Deliver a functional Android application and an iOS application from a single Kotlin Multiplatform project, sharing all business logic and as much user interface code as the platform allows, without sacrificing native look, feel, or performance.

\textbf{Objective 2: Adaptive AI-generated content with semantic deduplication.} Implement a nightly batch pipeline that generates domain-specific educational tasks using the Mistral language model API, validates their structure, and filters out semantically duplicate content using vector embeddings and cosine similarity search via pgvector, so that players always receive fresh content.

\textbf{Objective 3: Engagement loop engineering.} Implement a server-side scoring engine with configurable reward and penalty thresholds, badge awards, and streak tracking, such that the platform produces the behavioural reinforcement patterns identified by gamification research as drivers of sustained engagement.

\textbf{Objective 4: Cost efficiency through batch pre-generation.} Demonstrate that the nightly batch model produces a fixed daily cost independent of user count, in contrast to the on-demand model whose daily cost scales linearly with active users, and validate this model at the scale of the platform's initial deployment.

\section{Document Structure}

The remainder of this report is organised as follows. Chapter~\ref{chap:stateofart} surveys the relevant literature and existing technologies, covering gamification in education, serious games, multiplatform mobile development frameworks, AI content generation, and related platforms, and identifies the gap that this project addresses. Chapter~3 describes the system architecture in depth, following the engineering reasoning behind each design decision: the content generation pipeline, data persistence, mobile client design, authentication, the scoring engine, consistency trade-offs, observability, and deployment. Chapter~4 presents the implementation, detailing how the architectural decisions were realised in code, including the monorepo structure, Kotlin Multiplatform client, Spring Boot backend, LangChain4j pipeline, database schema, and Docker Compose deployment. Chapter~5 covers evaluation, reporting on functional testing, performance measurements, and a cost model validation. Chapter~6 draws conclusions, reflects on the engineering process, and outlines directions for future work.
